\chapter{Pendahuluan}

\section{Identitas Mata Kuliah}

Mata kuliah \textbf{IT30713 -- Enterprise \& Platform Architecture} merupakan
mata kuliah 3 SKS untuk mahasiswa magister di Pradita University. Mata kuliah
ini membahas cara \emph{Enterprise Architecture} (EA) membantu organisasi
menyelaraskan arah strategis dengan kapabilitas teknologi. Di dalamnya juga
dibahas peran strategi API dan tata kelola platform dalam membangun model
bisnis digital yang terhubung dengan ekosistem.

Mata kuliah ini dirancang untuk mahasiswa magister Indonesia yang sebagian
besar tidak berasal dari latar belakang IT. Karena itu, pembahasan diarahkan
pada cara berpikir strategis, manajerial, dan sistemik. Detail teknis tetap
diperkenalkan secukupnya, tetapi selalu ditempatkan sebagai alat bantu untuk
membuat keputusan, bukan sebagai tujuan utama pembelajaran.

\section{Capaian Pembelajaran}

\subsection{Capaian Pembelajaran Lulusan yang Didukung}

\begin{itemize}
    \item \textbf{CPL-4 -- System Design:} Mahasiswa menyusun artefak blueprint
    enterprise yang memetakan strategi, kapabilitas, data, aplikasi, teknologi,
    API, tata kelola, dan peta jalan sebagai dasar argumentasi makalah.
    \item \textbf{CPL-6 -- Leadership \& Collaboration:} Mahasiswa mengelola
    proyek penulisan individu, memberi dan menerima tinjauan sejawat, serta
    mengomunikasikan argumen rancangan kepada audiens teknis maupun non-teknis.
\end{itemize}

\subsection{Capaian Pembelajaran Mata Kuliah}

Setelah menyelesaikan mata kuliah ini, mahasiswa mampu:

\begin{enumerate}
    \item Menjelaskan peran Enterprise Architecture dalam transformasi digital
    organisasi.
    \item Menggunakan pendekatan TOGAF-lite untuk menyusun visi arsitektur,
    kondisi saat ini, arsitektur target, analisis kesenjangan, dan peta jalan.
    \item Memetakan domain Business, Data, Application, dan Technology (BDAT)
    sebagai dasar digital blueprint enterprise.
    \item Merancang strategi API untuk integrasi, kolaborasi, dan penciptaan
    nilai digital.
    \item Mengevaluasi model bisnis platform, network effects, ekosistem mitra,
    dan tata kelola platform.
    \item Merumuskan rekomendasi tata kelola, risiko, metrik, dan implikasi
    manajerial yang relevan untuk konteks Indonesia (UU PDP, UU ITE, OJK, BSSN,
    SNAP BI, PP 71/2019, Permenkominfo PSE).
    \item Menulis makalah ilmiah individual sesuai standar jurnal atau
    konferensi, termasuk struktur, sitasi, format, dan batas kemiripan naskah.
    \item Mempresentasikan makalah dalam format konferensi dan merespons umpan
    balik secara profesional.
\end{enumerate}

\section{Strategi Pembelajaran}

Pembelajaran menggunakan format \textbf{lokakarya individual}. Setiap sesi
berdurasi 3$\times$50 menit. Secara umum, 30 menit pertama digunakan untuk
pengantar konsep dan contoh kasus Indonesia, 90 menit berikutnya untuk
menyusun artefak atau bagian makalah secara individual, dan 30 menit terakhir
untuk konsultasi, tinjauan sejawat, serta refleksi kemajuan.

Setiap mahasiswa memilih satu organisasi studi kasus pada Sesi 1 dan
menggunakannya secara konsisten sampai akhir. Organisasi dapat berupa entitas
nyata atau hipotetis-realistis dalam konteks Indonesia, misalnya kampus, rumah
sakit, bank, koperasi atau BPR, retail, logistik UMKM, pemerintahan daerah,
agribisnis, atau BUMN. Istilah teknis dijelaskan melalui analogi manajerial dan
contoh organisasi.

\section{Luaran Utama dan Pelengkap}

\subsection{Luaran Utama}

Luaran utama mata kuliah adalah satu \textbf{makalah ilmiah individual} yang
diarahkan menjadi naskah siap diajukan setelah revisi final. Target publikasi
dapat berupa jurnal nasional terakreditasi (SINTA), jurnal internasional, atau
konferensi nasional/internasional yang relevan dengan Enterprise Architecture,
sistem informasi, transformasi digital, ekonomi API, atau strategi platform.

Seluruh aktivitas perkuliahan diarahkan untuk membangun bahan, kerangka,
argumen, artefak, dan kualitas akademik makalah secara bertahap, sesi demi
sesi.

\subsection{Luaran Pelengkap}

Luaran pelengkap wajib disusun sebagai bahan pembentuk makalah, tetapi tidak
memiliki bobot penilaian terpisah. Artefak ini berfungsi sebagai bukti analisis
dan bahan visual yang memperkuat makalah:

\begin{enumerate}
    \item Blueprint enterprise berbasis BDAT.
    \item Peta kapabilitas bisnis dan value stream.
    \item Visi arsitektur dan analisis pemangku kepentingan.
    \item Kondisi saat ini dan arsitektur target.
    \item Portofolio API atau peta kandidat API.
    \item Peta ekosistem platform atau kanvas model bisnis platform.
    \item Piagam tata kelola platform.
    \item Analisis kesenjangan dan peta jalan 12--24 bulan.
    \item Proposal makalah versi 1 (penanda kemajuan formatif pada Sesi 8).
    \item Bahan presentasi konferensi (Sesi 14).
    \item Surat pengantar atau daftar periksa pengajuan sesuai target publikasi.
\end{enumerate}

\section{Penilaian}

Penilaian akhir mata kuliah hanya berasal dari \textbf{makalah publikasi
individu} (100\%). Naskah harus dilengkapi artefak blueprint relevan sebagai
gambar, tabel, atau lampiran. Rincian rubrik disampaikan pada
Bab~\ref{ch:orientasi}. Rubrik tersebut menilai kejelasan masalah, kualitas
tinjauan pustaka, metode, kedalaman analisis, kualitas rancangan, diskusi dan
implikasi, kepekaan terhadap konteks Indonesia, kualitas akademik, serta
kelayakan publikasi.

Standar minimum naskah meliputi: mengikuti templat target publikasi, memiliki
minimal 15 referensi dengan setidaknya 5 referensi dari jurnal atau konferensi
terindeks, menggunakan gaya sitasi yang konsisten (APA, IEEE, atau gaya
target publikasi), menyertakan artefak blueprint yang benar-benar digunakan
dalam argumen, menjaga kemiripan Turnitin $\leq 25\%$ di luar daftar pustaka
dan kutipan wajar, serta mencantumkan keterangan penggunaan AI pada bagian
\emph{Acknowledgement} bila AI digunakan sebagai alat bantu.

\section{Rencana 14 Sesi Pembelajaran}

Tabel~\ref{tab:linimasa-sesi} menyajikan ikhtisar 14 sesi sekaligus linimasa
perkembangan makalah. Kolom \emph{Keluaran} menunjukkan hasil kerja mingguan
yang akan membentuk naskah final secara bertahap.

\begin{table}[htbp]
\centering
\scriptsize
\caption{Ikhtisar 14 sesi dan keluaran makalah per sesi.}
\label{tab:linimasa-sesi}
\begin{tabularx}{\textwidth}{@{}c X X@{}}
\hline
\textbf{Sesi} & \textbf{Topik} & \textbf{Keluaran untuk Makalah} \\
\hline
1  & Orientasi EA, Platform \& Penulisan Akademik & Pernyataan topik, kandidat target publikasi, motivasi awal \\
2  & Strategi Bisnis \& Rumusan Masalah & Rumusan masalah, pertanyaan penelitian (RQ), kontribusi awal \\
3  & TOGAF-lite, ADM \& Literature Review & Visi arsitektur, bibliografi awal ($\geq$10) \\
4  & Business Architecture & Peta kapabilitas, value stream, tinjauan pustaka v1 \\
5  & Data Architecture & Peta data konseptual, draf metode/pendekatan \\
6  & Application Architecture & Portofolio aplikasi, draf analisis \\
7  & Technology Architecture \& Kerangka Final & Arsitektur teknologi acuan, kerangka final \\
8  & Klinik Proposal Makalah (tengah semester) & \textbf{Proposal makalah v1} (penanda kemajuan formatif) \\
9  & API Strategy: API sebagai Produk & Draf analisis bagian strategi API \\
10 & API Design, Lifecycle \& Governance & Spesifikasi API, daftar periksa tata kelola \\
11 & Platform Economy \& Model Bisnis & Draf diskusi bagian platform \\
12 & Platform Governance, Etika \& Regulasi & Piagam tata kelola, draf diskusi dan kesimpulan \\
13 & Lokakarya Pra-Pengajuan & Draf final makalah, bahan presentasi, daftar periksa pengajuan \\
14 & Presentasi Konferensi \& Finalisasi & \textbf{Naskah final makalah ($100\%$)} \\
\hline
\end{tabularx}
\end{table}

\section{Aturan Tugas dan Etika Akademik}

\begin{enumerate}
    \item Seluruh tugas dikerjakan secara individu. Diskusi dan tinjauan
    sejawat diperbolehkan, tetapi naskah, analisis, dan kesimpulan harus
    merupakan karya masing-masing mahasiswa.
    \item Plagiarisme tidak diperkenankan. Kemiripan Turnitin $\leq 25\%$ di
    luar daftar pustaka dan kutipan wajar.
    \item AI boleh digunakan sebagai alat bantu untuk mencari ide awal,
    memperbaiki parafrase, atau memeriksa bahasa. Jika digunakan, mahasiswa
    wajib mencantumkan keterangannya pada bagian \emph{Acknowledgement}.
    Substansi analisis, argumen, dan kesimpulan tetap harus merupakan karya
    mahasiswa.
    \item Sitasi yang benar wajib untuk literatur, gambar, tabel, data, dan
    contoh kasus.
    \item Data sensitif organisasi dianonimkan atau digunakan dengan izin
    tertulis.
    \item Setiap rekomendasi dalam makalah memuat: alasan, dampak, risiko,
    prioritas, dan implikasi.
    \item Presentasi Sesi~14 wajib diikuti sebagai forum pertanggungjawaban
    akademik dan umpan balik final, namun nilai akhir tetap berasal dari
    makalah.
    \item Kehadiran minimal 75\%.
\end{enumerate}

\section{Catatan untuk Pembaca Mahasiswa}

Buku ini ditulis sebagai pendamping mata kuliah, bukan sebagai pengganti
literatur klasik di bidang Enterprise Architecture. Setiap bab dirancang agar
mahasiswa dapat membacanya secara mandiri sebelum pertemuan kelas, lalu
menggunakan sesi tatap muka untuk berlatih, berdiskusi, dan menerima umpan
balik.

Bahasa buku ini sengaja dibuat manajerial dan ringan. Jika menemukan istilah
teknis yang belum dikenal, pembaca dapat merujuk glosarium pada akhir setiap
bab atau membaca sumber tambahan yang tercantum pada bagian \emph{Bacaan
Lanjutan}.

Setiap bab diakhiri dengan latihan terapan individual yang menghasilkan
artefak atau narasi konkret. Artefak tersebut bukan tujuan akhir, melainkan
bahan penyusun yang akan dirangkai pada akhir semester menjadi satu makalah
publikasi yang utuh.
